\chapter*{Conferencia VI}
\addcontentsline{toc}{chapter}{LECTURA I - El timbre y el carácter del violín}
CABALLEROS: Desde nuestra última sesión, he completado los trabajos de reparación y retoque, y ahora les presento este viejo violín desgastado en condiciones de tocar. Como pueden observar, las marcas de la edad y el desgaste son mucho menos evidentes. Desde sus asientos, no se aprecian los parches en la cabeza y el clavijero. Las zonas desgastadas de la tapa armónica, el fondo y los aros se han recoloreado, y todo el exterior ha sido restaurado. De hecho, este viejo violín ya no parece una "cosa vieja sin valor". Pero el trabajo exterior, comparado con el interior, es insignificante. En el momento de la construcción de este violín, sus superficies exteriores recibieron una protección de barniz de calidad suficientemente duradera como para perdurar para siempre. La goma mástica, aunque tarda en secarse y nunca se endurece, es la goma más elástica y resistente que he observado. En mis manos, la goma mástica con aceite requirió un año para secarse; e incluso después de diez años, no se endurece. [Como saben, el barniz para violines es un tema que suscita interminables debates. No pretendo suscitar debate sobre este punto. Simplemente expondré los hechos tal como me parecen. Es un hecho, según mi observación, que cualquier barniz que seque y se endurezca resulta ser un daño grave para el sonido del violín. En mi experiencia, el barniz interfiere con la acción independiente de las fibras de la caja de resonancia exactamente en proporción a su rigidez una vez seco. Como experimento, muchas veces he "atado" el sonido aplicando gomas de secado rígido a la caja de resonancia. Igualmente, muchas veces lo he "desatado".
92
\label{cuidado_barniz}
tono mediante la eliminación de dicho barniz. No encuentro que la aplicación de gomas de secado rápido a la placa posterior y a los aros produzca ningún daño al tono. He conocido usuarios de violín que prefieren el
Tono "atado". Satisfacer el gusto de estos violinistas es fácil y seguro de lograr. Esta certeza se debe a que la acción del barniz, o más bien la falta de acción, es un factor constante. Dado que la acción de la madera no es un factor constante, la cantidad de barniz necesaria para "atar" el tono de cualquier caja de resonancia solo puede determinarse mediante prueba y error. Es un hecho que al tocar un violín con un tono "atado", se requiere menos cuidado al manejar el arco. En consecuencia, es más fácil hacerse pasar por un violinista hábil. Por supuesto, dicho tono no supera la prueba de larga distancia.
El barniz suave de este viejo violín, a pesar de su edad y uso, estaría hoy en perfectas condiciones si hubiera recibido el cuidado adecuado. Como ejemplo de cuidado correcto para el violín, no conozco a nadie que supere al renombrado violinista Bernhard Listemann. Que Listemann es un caballero de la vieja escuela no necesita más que una presentación. Que es un artista, el mundo lo sabe. Que es un conocedor y coleccionista de violines italianos antiguos, se evidencia en el momento en que abre su bóveda ignífuga. Al levantar uno de esos violines, su cuidado se evidencia así: Su mano izquierda sostiene el mástil; su pulgar e índice derechos sostienen el pasador de la cola. Al presentárselo, cortésmente
dice:	"Por favor, señor, sujételo así y no lo toque93

el barniz." En sus violines no se ve ni una mota de suciedad, ni una mota de resina, ni un rasguño de botón o uña, ni una huella grasienta de las yemas de los dedos. Tal es el cuidado adecuado del violín. Con tal cuidado continuado, las paredes de ladrillo y piedra de la mansión Listemann se convertirán en polvo antes de que el exterior de esos violines muestre signos de desintegración. Ahora bien, no estoy, en absoluto, delirando sobre el barniz de Cremona, ni sobre los colores de Cremona en barniz de Cremona. Nunca delira sin al menos un 50 por ciento de absoluto en el mío. Por lo tanto, cuando delira, es por una "buena" razón. Además, debido a peculiaridades en mis nervios ópticos (de las cuales no soy culpable), el "polvo aéreo" del período de Cremona no afecta mi vista. Por lo tanto, la goma copal transparente, templada con gomas transparentes de naturaleza más suave, me parece exactamente igual ya sea que esté sobre el Maggini, el Guarneri, el Amati, el Stradivari, el Montagnani, Francisco Ruggeri, o acostado sobre Franklino Robinsoni (Frank Robinson), pero, en lo que respecta al uso de colores, algunos de esos antiguos constructores de violines (no todos) poseían un desarrollo bastante inusual del sentido del color. Dado que el sentido del color suele ser una cuestión de desarrollo mediante la práctica, se deduce que este sentido no será poseído por todos en igual grado. Aunque todos los fabricantes antiguos podrían haber empleado las mismas gomas, sin embargo, al colorear dichas gomas, se podrían anticipar resultados muy diversos. Es evidente que la persona que coloreó el barniz aplicado a este viejo violín carecía en gran medida de sentido del color. Que su selección de goma fue
94

Su excelente calidad se evidencia en su resistencia al desgaste, a la vez que conserva su suavidad.
(Que alguien pudiera creer, ni por un instante, que las gomas pueden penetrar la madera sobre la que se aplican como barniz, me resulta asombroso. Incluso si la penetración por las gomas fuera posible, solo podría resultar en la ruina de la sonoridad. Tal ruina es fácil de demostrar. Muchas veces he arruinado la sonoridad de la caja de resonancia al aplicarle un material que sí penetra en la madera. Después de cada remojo, sea cual sea el agente, el tono permanece muerto. Me han traído varios violines arruinados por remojar la caja de resonancia. Para su tono muerto, no conozco otro remedio que una caja de resonancia nueva. ) El trabajo en las superficies internas de este viejo violín es de mucha mayor importancia para su tono. Esos delgados revestimientos y ligeros bloques de esquina se reemplazan por otros con más del doble de masa. Que la solidez de los revestimientos y los bloques de esquina aumenta la potencia del tono también es fácil de demostrar. Al retirar los forros y los bloques de las esquinas de cualquier violín con buena potencia tonal, la prueba así obtenida proporciona amplia evidencia. Al gran grosor original de esta caja de resonancia podemos atribuir el valor tonal restante de este violín. Si su grosor original se hubiera reducido al mismo grosor que el Stradivarius de 1707 (dada por Simoutre como 2 y 8-10, hasta 2 y 7-10 mm), esta caja de resonancia estaría hoy en un estado lamentable. La cantidad de desintegración por esas fuerzas destructivas, calor y humedad-
95

ture, ahora sería suficiente para la ruina total. Su regradación está copiada del Strad de 1707. Primero, les daré la oportunidad de juzgar su tono, y luego describiré la veta y el color dentro de la caja de resonancia. Como observan a corta distancia, la potencia de su tono no es grande; pero, sus otras cualidades tonales no pueden ser superadas. Observen particularmente la potencia de los armónicos. Como han observado, la potencia de los tonos armónicos, de diferentes violines, es una cantidad variable. De la madera densa de la caja de resonancia no he tenido éxito en obtener una gran potencia de tono armónico. Los armónicos son valiosos para el violín solista. De hecho, que los armónicos sobretonos y armónicos de un basset, como me gusta llamarlos, tonos resultantes, como los llaman los libros de texto, no puede haber un sonido más hermoso. Como se mostró anteriormente, a estos hermosos sonidos se debe atribuir el tono "rico" del violín. En mi experiencia, los tonos resultantes audibles son difíciles de producir. Su producción es tan difícil que, en muchos violines, no he podido producirlos en absoluto. Solo en cajas de resonancia que producen una pureza absoluta de sonido musical los he producido en grado notable. En todos los casos en que existen esas creaciones fantasmales, creo que el tono de tales violines ha alcanzado el límite del valor tonal. Para producir un tono resultante audible se requieren otros dos tonos en un acorde exacto. La exactitud en el acorde de los tonos genéricos es donde reside la dificultad. Ante la más mínima variación perceptible de la exactitud, esos sonidos filnay desaparecen instantáneamente, y ningún intento de persuasión puede
96

inducirlos a reaparecer hasta que se establezca la exactitud en los tonos genéricos. (Los libros de texto indican que, por regla general, los tonos resultantes tienen una altura dos octavas inferior a la del acorde genérico; sin embargo, también se mencionan excepciones. Por ejemplo: el tono resultante, producido por la tercera y la tónica, no está dos octavas por debajo de la tercera ni de la tónica, sino dos octavas por debajo de la quinta de esa escala en particular. Como ejemplo, trazo los tonos Si y su tónica Sol en el pentagrama; por lo tanto, la escala es Sol mayor y la quinta es Re. En este caso, el tono resultante está dos octavas por debajo de Re; y para hallar su altura, dividimos la altura de Re, 600 Hz, entre 4, lo que da como resultado 150 Hz. Los armónicos de Si y Sol están, respectivamente, una octava por encima. Así pues: Si equivale a 500 Hz; su armónico, a 1000 Hz; Sol equivale a 800 Hz; su armónico, a 1600 Hz. De nuevo, quiero llamar su atención sobre el hecho de que la combinación de estos cinco tonos de alturas tan diferentes es la causa del sonido rico del violín).
Sí, en verdad, he examinado cuidadosamente esas cajas de resonancia que producen el tono "rico". Incluso he vuelto a abrir algunos violines con el único propósito de reexaminar la caja de resonancia. Como saben, mi trabajo con el violín se ha centrado principalmente en aquellos que han estado en uso; por lo tanto, el "acabado" impide la determinación precisa de las cualidades físicas mediante el examen de las superficies exteriores. En la medida de mis posibilidades, ahora describiré dichas cualidades físicas y su color. Primero: La veta sigue una línea recta, sin desviación alguna, ni en ningún lugar. Segundo: Los crecimientos anuales no son ni los más anchos ni los más estrechos; y
97

medida en promedio, 18 por pulgada. (Los extremos son de 16 a 20.) Tercero: No hay apariencia de grasa alguna. Cuarto: No hay albura, ni mancha de madera negra. Quinto: La madera es quebradiza. Sexto: La veta entera es más bien blanda que densa. Séptimo: Se necesita una herramienta de corte afilada y suave para dejar una superficie lisa. Octavo: El color es el amarillo más profundo de la mantequilla con más rojo que en cualquier muestra de mantequilla. Noveno: Esta profundidad de color se extiende a través de la caja de resonancia.
Según todas las pruebas que he podido reunir en el ámbito de la carpintería, este tipo de cajas de resonancia solo se obtienen de los árboles más viejos y grandes.
Las tapas, los forros y los bloques de este antiguo violín están ahora sellados herméticamente. Por lo tanto, se espera que los agentes tan destructivos como el calor y la humedad nunca más vuelvan a afectarlo. Así pues, conservará su excelente sonido durante siglos. Ustedes mismos pueden comprobar la belleza de su sonido tras esta protección interna. También pueden comprobarlo con el hermoso sonido de otros 59 violines usados con una protección interna similar.
Alguien ha afirmado que la longevidad media de un violín usado es de 80 años. En este punto, la dificultad de obtener estadísticas precisas es evidente. Para el estudiante de violín, el período medio de utilidad del violín es un tema de gran interés. Como premisa, diré que el entusiasmo, más la fuerza del brazo que maneja el arco, son factores que influyen en la longevidad del violín y que no deben pasarse por alto. Por lo tanto, cuando digo que desgasté un buen violín...
98

Dentro de 30 años, lo primero que harás será mirarme a los ojos para ver si lo digo en serio o no. Después, sin duda, te fijarás en mis 90 kilos de peso y evaluarás mentalmente mis bíceps de 40 centímetros, además de mi entusiasmo. Ambos resistirán la prueba.
Me parece que, dado que mi violín no tiene nada de romántico, misterioso ni espectacular en su origen, su historia debe ser refrescante. Sin duda será única. Como importación, mi violín provino de Luxemburgo. En este hecho no hay nada extraordinario. Haller, bohemio él mismo, dijo que mi violín sin duda fue hecho por un pastor de ovejas en las montañas del Tirol. Yo digo bien por las ovejas, bien por la montaña, mejor para el pastor de las ovejas, y lo mejor de todo para mí, porque ese violín tenía un sonido. Por eso lo usé tanto que lo desgasté. Por eso Haller, mi profesor de violín, me pidió 100 dólares tres veces. Por eso no lo vendí. Por eso creo que no todos los buenos luthiers son de Cremona. Ese violín tenía un sonido potente. Pero su valor no radicaba en el sonido. Su intensidad tonal era excepcional. Su valor no estaba en la intensidad. Era brillante y dulce a la vez. Ahí no residía su valor. Otros violines poseían potencia, intensidad, brillantez y dulzura de tono, pero comparados con mi violín, no tenían valor para mí. Ahora he llegado a algo extraordinario. Mi violín poseía una calidad de tono similar a la humana. Su tono podía llorar en la tristeza más desesperada, en la contrición, en la alegría; podía orar con la más profunda devoción; podía reír con la máxima 99

podía dejarlo pasar; podía cantar dulcemente como un pájaro. En esas cualidades tonales reside su valor.
¿Vender mi violín?
¡Primero el hambre!
Incluso en ese caso, habría aprovechado al máximo las "ayudas económicas".
En treinta años de uso, ese tono poderoso, intenso, brillante, compasivo y humano se fue apagando en el piano. Al principio no descubrí la razón. La atribuí primero a una cosa y luego a otra. Probé cuerdas de diferentes tamaños; diapasones de diferentes alturas, de diferentes densidades, de diferentes pesos; puentes de diferente masa, de diferente densidad; postes, lo mismo; pero todo fue en vano. Hasta entonces, en momentos de triunfo, mi violín siempre me había mirado con brillo en busca de aprobación. Siempre le concedí todo lo que me pedía. Pero ahora noté un cambio en su mirada. En lugar de su brillo habitual, había tristeza. Me miraba con lástima. La profunda desesperación en su dulce rostro me iluminó. Al descubrir que las caricias de mi arco estaban aplastando la vida de mi mascota, el dolor en mi corazón era algo que quisiera olvidar. Suavemente lo coloqué en su estuche y no abrí ese estuche en dos
años. No pude.
Soy un violinista común y corriente. Nunca pude tocar nada detrás ni encima del puente, ni mucho menos delante de él. Con ese violín no tenía que hacer mucho. Funcionaba mejor cuando yo intervenía lo menos posible en sus estados de ánimo. A menudo cedía el paso a
100

esos estados de ánimo. Luego era un deleite. Luego me llevaba a una persecución a través de las sombras y el sol de la melodía de una manera que a la vez era mi desesperación por la representación mediante notas. Los años han pasado volando. Hoy la memoria me lleva de vuelta a los días en que vivía en la pequeña ciudad de
C. En su gran mayoría, su gente provenía de
hacia la costa atlántica. Cada uno de ellos tenía
experimentaron el dolor de las separaciones familiares. Cada uno de
tal, amorosa libertad, firme en propósito, impulsada por
esperanza, había dejado a sus seres queridos y se había puesto en camino hacia eso
brillante estrella occidental cuyo símbolo es conocido por el
El mundo como Estados Unidos Cada uno de ellos, sabiendo, el
crujidos, empujones, forcejeos, actividad humana, el
Dolor humano aplastante y abrumador en la partida, le pedí a mi violín que pintara esa escena abandonada por el pincel del artista por falta de color para el dolor del hermano
la separación del hermano, la separación del padre,
despidiéndose entre lágrimas, querida madre anciana; porque
El dolor de la hermana, sola, desafiando el mar en busca de un hogar.
con los libres; por el dolor de su padre y su familia
dejando, al sonido del gong, al tierno padre
tono en "adiós, hijo", "adiós, hija",
por su tono desgarrador hacia su esposa, "Madre" él
la llama, (gong), por el silencio como sus brazos
envuélvela, "Madre", es solo un susurro, se ha ido.
voz, (último gong), por el sollozo ahogado de
el pecho agitado por el dolor del padre, por su paso inestable,
mientras sigue ciegamente el sonido, por el golpe sordo de
"Madre" cayendo, por el silbido del vapor, por
''todos a bordo'', por el grito desmayado de desvanecerse
amigo, por el estruendo de las olas que golpean, por el
101

El grito enloquecedor del demonio de la tormenta, que llama a víctimas humanas que desafíen su furia, por la quietud de la calma que regresa, por el alivio de la tensión humana en ritmo 3-4, por el zumbido de la animación que regresa en un animado 2-4, por el grito alegre, "¡Tierra, ho!", sin embargo, mi violín aceptó tal invitación.
Me encantaba ese violín.
¡Pobre Patti! Ella, la única estrella brillante entre nosotros, los ancianos, a sus 65 años llorando desconsoladamente porque su otrora vibrante voz se ha ido apagando con la vejez.
En público no he llorado.
En privado, el pañuelo en mi derecha aún está
conveniente.
¡Qué extraño es cómo este ídolo de madera puede llegar a lo más profundo del corazón humano, desafiando desde allí a todo aquel que se le cruce!
Ocasionalmente, en momentos y lugares inesperados, uno lee declaraciones fragmentarias sobre tal o cual cosa capaz de restaurar el tono del violín. Después de que el tono de mi violín se hubiera deteriorado por completo, una de esas declaraciones llamó mi atención. El autor, sin firmar, afirmaba que el barniz de aceite, aplicado a violines desgastados, tenía el poder de restaurar el tono. En ese momento no pude ni afirmar ni negar tal afirmación porque no tenía experiencia con el barniz de aceite. Aunque esta afirmación carecía de lógica sólida, sin embargo, al ser algo que no había probado en mi violín, decidí darle una oportunidad al barniz de aceite. Le conseguí a un amigo una cantidad de dicho barniz, junto con instrucciones sobre cómo aplicarlo. Con las cuerdas y el puente en...
102

En esta etapa, comencé a aplicar barniz al óleo mediante el proceso de frotado. Como este método era nuevo para mí, no podía determinar la cantidad exacta que estaba usando, pero pensando que si hay algo bueno en poco, también lo hay en más, así que dediqué un día entero a frotarlo.
(En este momento, pienso en la verdad de: "Qué poco sabemos sin experiencia". También confieso que me ha dado un ataque de fiebre por el "regreso rápido". Ya saben lo que es la fiebre por el "regreso rápido". Es esa sensación de "no puedo esperar". La humanidad en general puede ser atacada por esta fiebre; pero su forma más virulenta se manifiesta en quienes tocan el violín). Por supuesto, al terminar el trabajo de hoy, estaba lleno de ansiedad, impulsado por la esperanza y abatido por la duda. La esperanza de poder restaurar el tono de mi violín destronado me hizo sentir demasiado ansioso por el regreso. Por lo tanto, tiré la almohadilla de frotar, tomé un arco y extraje del sol una octava del tono más "muerto".
imaginable.
(No muchos de nosotros pasamos por la vida sin ver fantasmas; al menos eso creíamos en ese momento, lo cual da para una historia. La cuestión es que los fantasmas aparecen de repente. Su llegada nunca se anuncia con un silbido estridente, ni con el tintineo de una campana, ni con el rechinido de una rueda, ni con el sonido de una trompeta. De repente vienen. De repente nos vamos; es decir, tan pronto como nuestro aliento, que se desvanece repentinamente, nos permite irnos. La repentina aparición fantasmal no es más sorprendente que el fenómeno que apareció en la caja de resonancia de mi violín cuando lo saqué de
103
.
debajo de mi barbilla. )
Este fenómeno consistió en una gran cantidad de aberturas, o protuberancias, con forma de cráter en el barniz de aceite que acababa de aplicar. Algunas de estas aberturas tenían un diámetro de 3 a 16 pulgadas, y se extendían desde ese diámetro hasta simples puntas. Las aberturas con forma de cráter estaban tan cerca unas de otras que el borde de cada una tocaba el de la contigua. En este fenómeno del barniz, lo más interesante es su ubicación. Sin duda, usted está familiarizado con la teoría del filósofo francés Savart sobre la cuestión de «Cómo funciona el violín para producir sonido». ¿Recuerda aquellos elaborados experimentos científicos y su conclusión: que la tapa, la tapa trasera y los aros se unen simultáneamente para golpear el aire contenido, produciendo así el sonido?
(Del libro de N. E. Simoutre, 1885, infiero que Savart todavía es considerado en Francia como una autoridad en esta cuestión. De otras fuentes, aprendo que las teorías de Savart se consideran refutadas. Hasta el momento de este fenómeno del barniz, fui seguidor de Savart. Sus borradores sobre ciencia parecían genuinos y, en ningún manual pude encontrar ni afirmación ni negación de su postura. Por lo tanto, acepté sus conclusiones sin intentar ninguna demostración por mí mismo. En realidad, después de que los filósofos alemanes e ingleses descartaran el tono del violín con sui generis (autogenerado),
Al abandonar prácticamente el problema por considerarlo irresoluble, no tenía ninguna esperanza de encontrar o ver jamás tal solución. No pretendo que dicha solución sea la solución para 104

existen ahora
excepto como solución parcial.
Desde mi punto de vista, la evidencia aportada
por este fenómeno, junto con otros
Las demostraciones, que se darán más adelante, son abundantes.
prueba de error en la teoría de Savart sobre cómo el violín
Funciona para producir sonido.
No cabe la menor duda sobre la causa de esta perturbación del barniz. La agitación de esta masa blanda se debe enteramente a oscilaciones o movimientos vibratorios (términos sinónimos) de la caja de resonancia. Apliqué el arco solo sobre la cuerda de sol. La cuestión es: ¿podemos usted, yo o cualquier otra persona predeterminar la ubicación de esta perturbación mediante alguna teoría existente sobre cómo funciona el violín para producir sonido? Es evidente que dicha perturbación debe existir en todo el cuerpo del violín, si todo el cuerpo actuara con la misma energía para producir sonido. El barniz que apliqué se distribuyó uniformemente por todo el cuerpo. El área de perturbación es limitada. En longitud, esta área es igual a
2 y pulgadas; su ancho, 1 y 1 pulgadas. Los cráteres más grandes se encuentran en el centro de dicha área. Desde
El centro, el tamaño de los cráteres disminuye hacia abajo.
a meros puntos. Es evidente que el más amplio
La oscilación de caja de resonancia se produjo directamente debajo de los cráteres más anchos. También es evidente por sí mismo.
que este punto es el mejor foro de debate
La oscilación es el centro de una determinada área responsable
para el tono de la cuerda G. Llamo la atención sobre el hecho de que
El tono G de este violín fue una vez conocido por su
poder. Llamo la atención sobre el hecho de que la ubicación de esta área de alteración del barniz ofrece una
105

valiosa pista en la regulación del tono del violín; también, al hecho de que el hallazgo del área responsable del tono de la cuerda G condujo al hallazgo de las áreas responsables del tono de las cuerdas D, A y E; también, a mi conclusión de que la caja de resonancia del violín es totalmente responsable del tono del violín. Me refiero a que a los golpes entregados por la caja de resonancia sobre contenido
¿Qué características del sonido del violín atribuyo ahora? (Esta afirmación no incluye modificadores del sonido del violín, como la tapa trasera fuerte; la tapa trasera débil; la superficie interna imperfecta de la tapa trasera; la posición y el área de las salidas; el modelo del violín; la posición, longitud, densidad y diámetro del poste; la posición, densidad, altura y masa del puente; la altura, el peso y la superficie inferior del diapasón; el diámetro y la calidad de las cuerdas; la cantidad y la calidad del barniz. Estos modificadores del sonido se considerarán más adelante).
Ahora analizaré la ubicación de la alteración del barniz. El centro de esta zona se encuentra a medio camino entre el puente y el borde superior de la caja de resonancia, ligeramente a la izquierda de la barra. Probablemente nadie verá la ubicación de esta zona de alteración del barniz como yo la veo. Desde mi perspectiva, este fenómeno del barniz, y su ubicación, ilumina directamente las zonas de la caja de resonancia responsables del timbre de cada cuerda del violín.
Durante más de diez años de trabajo continuo, desde que ocurrió este accidente (lo llamo accidente n.° 2), he utilizado este índice en la regulación del tono de cajas de resonancia usadas sin encontrarme nunca con una decepción. Por lo tanto, estoy convencido de que este
106
\label{plancha_tono}
El índice es fiable. Gracias a él aprendí dónde encontrar y cómo corregir errores en la caja de resonancia que perjudican el tono de todas las cuerdas, o solo de una en particular. Por lo tanto, considero que este índice es de suma importancia para el violín. También reconozco que este índice es fruto del azar. La ciencia no tiene nada que ver con su existencia. La ciencia solo puede ofrecer una explicación a posteriori. Sigo creyendo que la ciencia no puede construir un violín. Es decir, que ninguna fórmula científica puede, ni tú, ni yo, ni nadie, predeterminar el tono del violín en función de todas sus peculiaridades.
Quiero destacar que este índice simplifica enormemente la regulación del tono; asimismo, que dicha regulación puede dirigirse exclusivamente a la caja de resonancia. En este sentido, la evidencia aportada por el accidente n.° 2 puede aceptarse como corroboración de los hechos conocidos desde hace años por los estudiantes de violín. En los últimos diez años de mi trabajo, he ignorado por completo la tapa trasera como elemento generador de tono. Durante ese tiempo, solo la he considerado como un modificador del tono. Ahora, lo único que le pido a la tapa trasera es que su superficie interior sea absolutamente perfecta para la reflexión de las ondas sonoras y que su rigidez sea suficiente para soportar, sin temblor, la carga del movimiento molecular que se origina en la caja de resonancia. Soy consciente de que muchos buenos luthiers siguen considerando la tapa trasera como un elemento generador de tono. Dichos luthiers pueden seguir produciendo buenos violines.
107

(Esta cuestión parece tener una tenacidad tan grande como la del cultivo de patatas. El emprendedor circasiano, si bien descubrió que los aborígenes americanos disfrutaban de las patatas, no encontró ninguna autoridad tradicional que indicara la fase lunar adecuada para su cultivo. Esta omisión provocó una división aparentemente perpetua entre los circasianos. Sin embargo, curiosamente, ambos bandos en esta controversia cultivan buenas patatas.)
[*]Al observar la ubicación de esta alteración del barniz, también se demuestra el poder de la acción simpática. En resumen, como ley física, se puede afirmar que el poder de la acción simpática disminuye inversamente con el cuadrado de la distancia. Por lo tanto, la parte de la caja de resonancia más cercana a la cuerda debe recibir mayor impulso de la acción de la cuerda que la parte más alejada de las cuerdas. Por consiguiente, podríamos presuponer que la mayor alteración del barniz se encuentra en la mitad superior de la caja de resonancia. Al sostener el violín bajo mi barbilla, pude sentir claramente la vibración de la caja de resonancia en ese punto, como era habitual. Por lo tanto, sé que la caja de resonancia actuó en casi toda su longitud en el momento de causar la alteración del barniz. Pero que dicha acción fue de menor potencia en la mitad inferior de la caja de resonancia queda demostrado por el hecho de que no apareció ninguna alteración del barniz en la mitad inferior.
(Dentro de los límites de mi observación, acción simpática, excitada en la caja de resonancia por
acción de cuerda, no ha recibido hasta ahora atención-
108

Por eso, deseo conocer su opinión al respecto. De hecho, al dar mis conclusiones sobre este tema, espero contar con el respaldo de otros estudiosos del sonido del violín. Dado que esta nueva cuestión surge de forma fortuita, no siento ningún apego a ella. Es tanto suya como mía. Solo la formulación de la pregunta me pertenece por casualidad.
No considero complejas las leyes que rigen la acción solidaria. Al contrario, son fáciles de comprender. Hasta el momento, solo puedo mencionar dos hechos como fundamentos de dichas leyes:
(1) Igual susceptibilidad a la fuerza. (2) Proximidad de los cuerpos.
Que la igual susceptibilidad a la fuerza es una condición necesaria para la acción simpática puede demostrarse intentando excitar dicha acción en cuerpos que poseen masa, densidad y rigidez muy variables. Tales intentos fracasan porque cuerpos tan variados no se excitan a la acción vibratoria con una fuerza idéntica. Como medio a mano para demostrar la susceptibilidad desigual a la fuerza, empleo este violín. Como pueden observar, cada cuerda de este violín produce un tono potente. Estas cuerdas son de calibre 2 cuidadosamente seleccionadas. La rigidez de la caja de resonancia debajo de cada cuerda se ha reducido cuidadosamente hasta que la fuerza en el golpe de cada cuerda, a tono de concierto, es suficiente para producir acción vibratoria en la caja de resonancia y corresponde a la acción de cada cuerda. Quito este G y lo reemplazo con este E. En su nueva posición, el tono E es débil. Ahora, el sonido-
109

La caja de resonancia debajo del Mi no está sujeta a la misma fuerza que el Mi. Por lo tanto, no existe interacción simpática entre estos dos cuerpos. En su posición correcta, el sonido de este Mi es potente; y dicha potencia sonora se debe a que el Mi y la caja de resonancia debajo del Mi están sujetos a la misma fuerza. Por consiguiente, existe interacción simpática, y su existencia aumenta la potencia sonora.
(A menudo, en trabajos de retonación, demostré el valor de "Dejar las cosas como están". Así, después de determinar la longitud de las cuerdas a utilizar, después de determinar la longitud por la posición del puente, después de ajustar la longitud, la masa y la posición del poste, después de determinar la masa y la altura del puente, después de probar el tono y encontrar pequeñas debilidades aquí y allá, y sabiendo con certeza que dicha debilidad se debía a demasiada madera de la caja de resonancia aquí y allá, después de reducir el grosor en tales lugares con la máxima precaución, después de asegurar una potencia uniforme para todas las cuerdas, entonces, como un idiota, en lugar de dejar las cosas como están, he perdido mi trabajo por intentar
hacerlo mejor que "suficientemente bien". El ' 'Elgin,"
ajustado a la temperatura, posición, isocronismo y funcionando "suficientemente bien", no es más susceptible a un trato idiota que el violín finamente ajustado. Con la misma seguridad no podemos limar metal en ningún punto de la periferia del volante de Elgin, que no podemos cambiar el tono de las cuerdas, el diámetro de las cuerdas, la longitud de las cuerdas, la altura del puente, la masa del puente, la densidad del puente, la longitud del poste, la masa del poste, la posición del poste y la rigidez de la caja de resonancia.
110

La calidad del violín... finamente ajustado y con el tono regulado. "No sirve de nada llorar por arruinar mi trabajo.Llorar solo añade "bebé" a "idiota". Con expresiones más o menos coloreadas, me levanto de un salto, agarro el lápiz rojo y escribo en la pared: "Que sea suficiente".
Solo.")
El familiar y divertido violín de palo de escoba no carece de una valiosa lección. Aunque su única cuerda vipjin D no despierta ninguna "simpatía" ni en el público ni en el palo de escoba, tiene la oportunidad de demostrar su poder de producción de tono cuando no recibe ayuda alguna de la simpatía.
acción.	
La proximidad de los cuerpos, como base para la ley que rige la acción simpática, puede demostrarse de diversas maneras. Como método imperfecto, y solo por conveniencia, alejo todas las cuerdas de este violín de la caja de resonancia, reemplazando el puente por otro una pulgada más alto. Como pueden observar, la potencia sonora de todas las cuerdas disminuye perceptiblemente. Si bien es imperfecta, esta demostración muestra que la acción simpática disminuye a medida que aumenta la distancia.
Es evidente que la distancia puede aniquilar la acción simpática. Por el contrario, la proximidad aumenta la acción simpática. De estos hechos concluyo que la acción simpática causa una mayor amplitud de oscilación en la mitad superior de la caja de resonancia; y, a dicha mayor oscilación, con su potencia aumentada, atribuyo la ubicación de esta perturbación del barniz. De los mismos hechos, también concluyo que la mitad superior de la caja de resonancia,
111

a una distancia limitada a cada lado de la barra, es responsable de la potencia tonal de la cuerda G. De los mismos hechos concluyo que la rigidez debería ser menor en la mitad inferior que en la mitad superior de la caja de resonancia. La distancia ciertamente disminuye la fuerza de la acción simpática como el cuadrado de la distancia. (Sabiendo que cada frase en la descripción de este fenómeno del barniz, junto con mis conclusiones de la misma, estará sujeta a los límites del escrutinio, por lo tanto he recurrido a mi máxima capacidad de claridad de dicción. Si encuentra falta de claridad en algún punto, y si desea una mayor aclaración al respecto, solo necesita notificármelo. De antemano, solicito sus opiniones sobre mis conclusiones. Admito el hecho de que "dos cabezas piensan mejor que una". Admito que
Puede que tus opiniones sean mejores que las mías; por lo tanto, serán recibidas con agrado.
Ahora he completado la presentación del fenómeno del barniz n.º 1. Fenómeno del barniz n.º 1.
El número 2 se presentará en una ocasión posterior.
